% GENERATED FILE. Edit canonical lessons, then run npm run generate:books.
% canonical-source-hash: fnv1a64:b790b59535426ac1
% canonical-lessons: HI-C205-patra, HI-C205-lifafa, HI-C205-tikatghar, HI-C205-sandesa, HI-C205-batcit

\chapter{A letter, An envelope, A ticket office, A message, A conversation}
\label{ch:hi-a-letter-an-envelope-a-ticket-office-a-message-a-conversation}
\hlchaptermodality{205}

\begin{chapteropening}
\textbf{By the end of this chapter:} \emph{I can say a letter, an envelope, a ticket office, a message and a conversation.}

Complete the last lesson of chapter 205: I can say a letter, an envelope, a ticket office, a message and a conversation.
\end{chapteropening}

\section[\texorpdfstring{\hi{पत्र} (patra)}{पत्र (patra)}]{\hi{पत्र} (patra) — a letter (formal)}

\label{lesson:HI-C205-patra}



\begin{quote}
Before the new one: say the Hindi for petrol, then the Hindi for a telephone.
\end{quote}

\subsection*{You'll want to know: \hi{पत्र}}
\textbf{\hi{पत्र}} — \emph{patra} — \textquotedblleft{}a letter (formal)\textquotedblright{}.

\begin{cousinweb}
\textbf{\hi{पत्र}} (\emph{patra}) is Sanskrit \emph{pattra}, \textquotedblleft{}wing-feather; leaf; leaf of a book\textquotedblright{}, taken over whole as a learned loan, a tatsama (CDIAL 7733). A letter is a leaf of writing. CDIAL lists the inherited form as \emph{pattā}, \textquotedblleft{}leaf\textquotedblright{}.
\end{cousinweb}

\begin{grammarlens}[title={What you've built}]
One more word for this chapter.
\end{grammarlens}

\subsection*{Your turn}
\begin{itemize}
\raggedright
  \item \emph{Say it:} \emph{patra}
  \item \emph{Say it:} \emph{patra}, once more
  \item \emph{Say it:} say \emph{patra}
  \item \emph{Recall:} say the Hindi for petrol, then the Hindi for a telephone, then say \emph{patra} again
\end{itemize}

\begin{tcolorbox}[breakable,colback=teal!4,colframe=teal!35!black,title={Before you move on}]
What does \hi{पत्र} mean? (\textquotedblleft{}A letter (formal)\textquotedblright{}.) Say it once more.
\end{tcolorbox}
\section[\texorpdfstring{\hi{लिफ़ाफ़ा} (lifāfā)}{लिफ़ाफ़ा (lifāfā)}]{\hi{लिफ़ाफ़ा} (lifāfā) — an envelope}

\label{lesson:HI-C205-lifafa}



\begin{quote}
Before the new one: say the Hindi for a telephone, then the Hindi for a letter.
\end{quote}

\subsection*{You'll want to know: \hi{लिफ़ाफ़ा}}
\textbf{\hi{लिफ़ाफ़ा}} — \emph{lifāfā} — \textquotedblleft{}an envelope\textquotedblright{}.

\begin{cousinweb}
\textbf{\hi{लिफ़ाफ़ा}} (\emph{lifāfā}) is a loan from Arabic \emph{lifāfa}, \textquotedblleft{}a wrapper\textquotedblright{}, which reached Hindi through Persian. Its Arabic root is \emph{l-f-f} (Lane).
\end{cousinweb}

\begin{grammarlens}[title={What you've built}]
One more word for this chapter.
\end{grammarlens}

\subsection*{Your turn}
\begin{itemize}
\raggedright
  \item \emph{Say it:} \emph{lifāfā}
  \item \emph{Say it:} \emph{lifāfā}, once more
  \item \emph{Say it:} say \emph{lifāfā}
  \item \emph{Recall:} say the Hindi for a telephone, then the Hindi for a letter, then say \emph{lifāfā} again
\end{itemize}

\begin{tcolorbox}[breakable,colback=teal!4,colframe=teal!35!black,title={Before you move on}]
What does \hi{लिफ़ाफ़ा} mean? (\textquotedblleft{}An envelope\textquotedblright{}.) Say it once more.
\end{tcolorbox}
\section[\texorpdfstring{\hi{टिकट} \hi{घर} (ṭikaṭ ghar)}{टिकट घर (ṭikaṭ ghar)}]{\hi{टिकट} \hi{घर} (ṭikaṭ ghar) — a ticket office}

\label{lesson:HI-C205-tikatghar}



\begin{quote}
Before the new one: say the Hindi for a letter, then the Hindi for an envelope.
\end{quote}

\subsection*{You'll want to know: \hi{टिकट} \hi{घर}}
\textbf{\hi{टिकट} \hi{घर}} — \emph{ṭikaṭ ghar} — \textquotedblleft{}a ticket office\textquotedblright{}.

\begin{cousinweb}
\textbf{\hi{टिकट} \hi{घर}} (\emph{ṭikaṭ ghar}) joins English \emph{ticket} with \emph{ghar}, \textquotedblleft{}house\textquotedblright{}, from \emph{ghara} (CDIAL 4428).
\end{cousinweb}

\begin{grammarlens}[title={What you've built}]
One more word for this chapter.
\end{grammarlens}

\subsection*{Your turn}
\begin{itemize}
\raggedright
  \item \emph{Say it:} \emph{ṭikaṭ ghar}
  \item \emph{Say it:} \emph{ṭikaṭ ghar}, once more
  \item \emph{Say it:} say \emph{ṭikaṭ ghar}
  \item \emph{Recall:} say the Hindi for a letter, then the Hindi for an envelope, then say \emph{ṭikaṭ ghar} again
\end{itemize}

\begin{tcolorbox}[breakable,colback=teal!4,colframe=teal!35!black,title={Before you move on}]
What does \hi{टिकट} \hi{घर} mean? (\textquotedblleft{}A ticket office\textquotedblright{}.) Say it once more.
\end{tcolorbox}
\section[\texorpdfstring{\hi{संदेसा} (sandesā)}{संदेसा (sandesā)}]{\hi{संदेसा} (sandesā) — a message (everyday)}

\label{lesson:HI-C205-sandesa}



\begin{quote}
Before the new one: say the Hindi for an envelope, then the Hindi for a ticket office.
\end{quote}

\subsection*{You'll want to know: \hi{संदेसा}}
\textbf{\hi{संदेसा}} — \emph{sandesā} — \textquotedblleft{}a message (everyday)\textquotedblright{}.

\begin{cousinweb}
\textbf{\hi{संदेसा}} (\emph{sandesā}) is a tadbhava, a word inherited from Sanskrit and reshaped by sound change. It comes from \emph{saṁdeśa}, \textquotedblleft{}message, information\textquotedblright{} (CDIAL 12904). It is the inherited form of the word Hindi also has whole as \emph{sandeś}.
\end{cousinweb}

\begin{grammarlens}[title={What you've built}]
One more word for this chapter.
\end{grammarlens}

\subsection*{Your turn}
\begin{itemize}
\raggedright
  \item \emph{Say it:} \emph{sandesā}
  \item \emph{Say it:} \emph{sandesā}, once more
  \item \emph{Say it:} say \emph{sandesā}
  \item \emph{Recall:} say the Hindi for an envelope, then the Hindi for a ticket office, then say \emph{sandesā} again
\end{itemize}

\begin{tcolorbox}[breakable,colback=teal!4,colframe=teal!35!black,title={Before you move on}]
What does \hi{संदेसा} mean? (\textquotedblleft{}A message (everyday)\textquotedblright{}.) Say it once more.
\end{tcolorbox}
\section[\texorpdfstring{\hi{बातचीत} (bātcīt)}{बातचीत (bātcīt)}]{\hi{बातचीत} (bātcīt) — a conversation}

\label{lesson:HI-C205-batcit}



\begin{quote}
Before the new one: say the Hindi for a ticket office, then the Hindi for a message.
\end{quote}

\subsection*{You'll want to know: \hi{बातचीत}}
\textbf{\hi{बातचीत}} — \emph{bātcīt} — \textquotedblleft{}a conversation\textquotedblright{}.

\begin{cousinweb}
\textbf{\hi{बातचीत}} (\emph{bātcīt}) is a rhyming pair: \emph{bāt}, \textquotedblleft{}talk\textquotedblright{}, from Sanskrit \emph{vārttā} (CDIAL 11564), and \emph{cīt}, which CDIAL lists in this very pair under \emph{citti}, \textquotedblleft{}thought\textquotedblright{} (CDIAL 4801).
\end{cousinweb}

\begin{grammarlens}[title={What you've built}]
One more word for this chapter.
\end{grammarlens}

\subsection*{Your turn}
\begin{itemize}
\raggedright
  \item \emph{Say it:} \emph{bātcīt}
  \item \emph{Say it:} \emph{bātcīt}, once more
  \item \emph{Say it:} say \emph{bātcīt}
  \item \emph{Recall:} say the Hindi for a ticket office, then the Hindi for a message, then say \emph{bātcīt} again
\end{itemize}

\begin{tcolorbox}[breakable,colback=teal!4,colframe=teal!35!black,title={Before you move on}]
What does \hi{बातचीत} mean? (\textquotedblleft{}A conversation\textquotedblright{}.) Say it once more.
\end{tcolorbox}
